\subsubsection{Page-Fault Error Code}
\manualtablecaption{PAGE\_FAULT ERROR\_CODE Fields}
\begin{manuallongtable}{@{}>{\raggedright\arraybackslash}p{0.70in}>{\raggedright\arraybackslash}p{1.45in}>{\raggedright\arraybackslash}p{3.35in}@{}}
\toprule
\rowcolor{ManualHeaderFill}
\textbf{Bits} & \textbf{Field} & \textbf{Meaning}\\
\midrule
\endfirsthead
\multicolumn{3}{l}{\scriptsize\itshape Table \themanualtable\ (continued)}\\
\toprule
\rowcolor{ManualHeaderFill}
\textbf{Bits} & \textbf{Field} & \textbf{Meaning}\\
\midrule
\endhead
7..0 & \texttt{CAUSE} & page-fault cause listed below\\
9..8 & \texttt{ACCESS} & 0 none, 1 read, 2 write, 3 execute\\
10 & \texttt{USER\_DOMAIN} & access selected the user-domain permission path\\
11 & \texttt{ATOMIC} & access was an atomic read-modify-write\\
14..12 & \texttt{WALK\_LEVEL} & 0 before or without a walk; 1 through 5 identify the PTE level\\
15 & reserved & zero\\
23..16 & \texttt{OPERAND} & zero-based explicit operand ordinal; \texttt{0xff} for implicit fetch or stack access\\
24 & \texttt{FAULT\_EA\_VALID} & \texttt{FAULT\_EA} was produced\\
25 & \texttt{FAULT\_LINEAR\_VALID} & \texttt{FAULT\_LINEAR} was produced\\
26 & reserved & zero\\
29..27 & \texttt{ACCESS\_SIZE} & 0 unavailable or inapplicable; 1 B; 2 W; 3 L; 4 Q\\
63..30 & reserved & zero\\
\bottomrule
\end{manuallongtable}

\manualtablecaption{PAGE\_FAULT Causes}
\begin{manualtabular}{@{}>{\raggedright\arraybackslash}p{0.65in}>{\raggedright\arraybackslash}p{2.20in}>{\raggedright\arraybackslash}p{2.55in}@{}}
\toprule
\rowcolor{ManualHeaderFill}
\textbf{Value} & \textbf{Cause} & \textbf{Detection class}\\
\midrule
0 & \texttt{NOT\_PRESENT} & required PTE is not present\\
1 & \texttt{PERMISSION} & accumulated permission excludes the access\\
2 & \texttt{MALFORMED\_PTE} & structurally invalid PTE\\
3 & \texttt{NONCANONICAL} & noncanonical address\\
4 & \texttt{SEGMENT\_BOUNDS} & segment bounds failure\\
5 & \texttt{ATOMIC\_ALIGNMENT} & misaligned byte-addressed (\texttt{AT=0}) atomic operand\\
6 & \texttt{ADDRESS\_TYPE} & address or memory type rejects the access\\
\bottomrule
\end{manualtabular}

The originating privilege is recovered from saved \texttt{STATUS.PM}; it is not duplicated in \texttt{ERROR\_CODE}.
\texttt{USER\_DOMAIN} instead distinguishes the ordinary user permission path from the checked user-access path used
by supervisor instructions such as \texttt{MOVUC}, \texttt{MOVCU}, and \texttt{MOVUU}. An atomic read-modify-write
reports \texttt{ACCESS=write}. \texttt{NOT\_PRESENT} and \texttt{MALFORMED\_PTE} identify the actual PTE level.
\texttt{PERMISSION} identifies the first level, in root-to-leaf walk order, whose accumulated permissions exclude the
requested access. A cause detected before a walk reports level zero.

\texttt{ACCESS\_SIZE} records the architectural transfer width when one of the B, W, L, or Q widths applies. It is
zero when no such width is available or applicable.

\texttt{FAULT\_EA} is the numeric effective address before segment pre-translation. \texttt{FAULT\_LINEAR} is the
result after segment pre-translation. A value not produced is zero and has its corresponding validity bit clear.

\subsubsection{Other Exception Error Codes}
\texttt{DIVIDE\_ERROR.ERROR\_CODE[7:0]} is 0 for \texttt{ZERO\_DIVISOR} and 1 for
\texttt{SIGNED\_OVERFLOW}; bits 63..8 are zero.

\manualtablecaption{ILLEGAL\_INSTRUCTION Causes}
\begin{manualtabular}{@{}>{\raggedright\arraybackslash}p{0.65in}>{\raggedright\arraybackslash}p{2.35in}>{\raggedright\arraybackslash}p{2.40in}@{}}
\toprule
\rowcolor{ManualHeaderFill}
\textbf{Value} & \textbf{Cause} & \textbf{Meaning}\\
\midrule
0 & \texttt{INVALID\_OPCODE} & opcode is not assigned\\
1 & \texttt{INVALID\_EA\_FORM} & effective-address form is not legal for the instruction\\
2 & \texttt{RESERVED\_ENCODING} & assigned instruction uses a reserved field value\\
3 & \texttt{UNAVAILABLE\_EXTENSION} & required extension is unavailable\\
4 & \texttt{EXPLICIT\_ILLEGAL} & architected \texttt{ILLEGAL} instruction\\
\bottomrule
\end{manualtabular}

\manualtablecaption{INVALID\_CONTROL\_STATE Causes}
\begin{manualtabular}{@{}>{\raggedright\arraybackslash}p{0.65in}>{\raggedright\arraybackslash}p{2.35in}>{\raggedright\arraybackslash}p{2.40in}@{}}
\toprule
\rowcolor{ManualHeaderFill}
\textbf{Value} & \textbf{Cause} & \textbf{Meaning}\\
\midrule
0 & \texttt{INVALID\_SELECTOR} & selector names no available resource\\
1 & \texttt{RESERVED\_BITS} & reserved bits are nonzero or changed illegally\\
2 & \texttt{INVALID\_IMAGE} & supplied architectural image is invalid\\
3 & \texttt{INVALID\_TRANSITION} & requested state transition is illegal\\
\bottomrule
\end{manualtabular}

An unavailable \texttt{RDPMC} ID reports \texttt{INVALID\_SELECTOR}; an illegal
reserved-bit update reports \texttt{RESERVED\_BITS}; an invalid segment,
control-register, or SAVE/RESTORE image reports \texttt{INVALID\_IMAGE}; and an
illegal \texttt{ERET}, \texttt{SYSRET}, or privilege-state transition reports
\texttt{INVALID\_TRANSITION}.

\texttt{FLOATING\_POINT\_FAULT.ERROR\_CODE[4:0]} is a cause bitmap: bits 0 through 4 are \texttt{NV},
\texttt{DZ}, \texttt{OF}, \texttt{UF}, and \texttt{NX}. All causes generated by the operation are reported together;
bits 63..5 are zero.

The bounds instructions report failure through \texttt{FLAGS.V} and do not raise \texttt{BOUNDS\_FAULT}; segment
bounds failures use \texttt{PAGE\_FAULT}. This specification also has no \texttt{ARITHMETIC\_TRAP} producer. Both exception
IDs remain reserved. \texttt{BKPT} raises \texttt{BREAKPOINT}; \texttt{DEBUG\_TRACE} is generated only by the
architectural trace mechanism.
